% Part: turing-machines
% Chapter: undecidability
% Section: trakhtenbrot

\documentclass[../../../include/open-logic-section]{subfiles}

\begin{document}

\olfileid{tur}{und}{tra}
\olsection{Teorema Trakhtenbrot}

\begin{explain}
  Ing \olref[rep]{sec} kita netepake !!{sentence}s
  $!T(M,w)$ lan~$!E(M,w)$ kanggo mesin Turing~$M$
  lan rerangken input~$w$. Banjur ing
  \olref[ver]{lem:valid-if-halt} lan
  \olref[ver]{lem:halt-if-valid} kita nuduhake yen
  $!T(M,w) \lif !E(M,w)$ sah yen lan mung yen $M$,
  nalika diwiwiti nganggo input~$w$, pungkasane mandheg. Cathetan OLPL-773 lan SOL6-F058: klausa pangreksan lan bukti representasi sadurunge saiki dibenerake supaya kothak asal kang ditulisi ora dikekang nganggo simbol lawas.
  Amarga Masalah Mandheg ora bisa diputusake, iki
  ateges kasahihan lan panyukupan !!{sentence}s
  logika tataran kapisan ora bisa diputusake
  (\cref{tur:und:uns:thm:decision-prob,tur:und:uns:cor:undecidable-sat}).

  Nanging kasahihan lan panyukupan ukara ditegesi
  tumrap sembarang !!{structure}s, kang winates utawa
  tanpa wates. Bisa uga kok kira luwih gampang
  mutusake apa !!a{sentence} bisa dicukupi ing
  !!{structure} winates (utawa sah ing kabeh
  !!{structure}s winates). Bukti yen masalah putusan
  ora bisa dirampungake bisa kita adaptasi kanggo
  nuduhake yen prakara iki uga ora luwih gampang.

  Kapisan, yen bali menyang bukti
  \olref[ver]{lem:halt-if-valid}, bakal katon yen ing kana
  kita mbangun model~$\Struct M$ saka~$!T(M,w)$ kang
  njlentrehake persis tumindake mesin~$M$ nalika
  diwiwiti nganggo input~$w$. Domaine model iku
  yaiku~$\Nat$, dadi tanpa wates. Nanging yen $M$
  pancen mandheg nganggo input~$w$, kita bisa
  mbangun model winates~$\Struct M'$ kanthi cara
  kang padha. Upamane $M$ kang diwiwiti nganggo
  input~$w$ mandheg sawisé~$k$ langkah. Dadekna
  domain~$\Domain{M'}$ minangka himpunan
  $\{0, \dots, n\}$, ing ngendi $n$ iku nilai
  kang luwih gedhé tinimbang~$k$ kang ditambah siji lan dawa~$w$,
  banjur tetepna
  \[
    \Assign{\prime}{M'}(x) = 
    \begin{cases}
      x + 1 &\text{yen $x < n$}\\
      n &\text{yen ora,}
    \end{cases}
  \]
  sarta $\tuple{x,y} \in \Assign{<}{M'}$ yen lan
  mung yen $x < y$ utawa $x = y = n$. Kanggo
  bagean liyane, $\Struct{M'}$ ditegesi kaya~$\Struct{M}$.
  Miturut definisi~$\Struct{M'}$, kaya ing bukti
  \olref[ver]{lem:halt-if-valid}, $\Sat{M'}{!T(M,w)}$. Cathetan OLPL-161 lan SOL6-F058: pilih pucuk domain ngluwihi kabeh wektu transisi, kothak kang ditekani lan input. Klausa pangreksan saiki bener; konfigurasi mandheg pungkasan digunakake maneh kanggo wektu sabanjure.
  Lan amarga kita nganggep $M$ mandheg nganggo
  input~$w$, $\Sat{M'}{!E(M,w)}$. Dadi,
  $\Struct{M'}$ iku model winates saka
  $!T(M,w) \land !E(M,w)$ (gatekna yen $\lif$
  wis diganti dadi~$\land$).
  
  Iki wis separo dalan menyang bukti: kita wis
  nuduhake yen $M$ mandheg nganggo input~$w$,
  mula $!T(M,w) \land !E(M,w)$ nduwé model winates.
  Sayange, kosok baline ora bener: ana mesin Turing
  kang ora mandheg kanggo input~$w$ tartamtu,
  nanging $!T(M,w) \land !E(M,w)$ tetep nduwé
  model winates. Tuladhane, gatekna mesin~$M$
  kang mung nduwé kahanan $q_0$ lan instruksi
  $\delta(q_0,\TMblank) = \tuple{q_0,\TMblank,\TMstay}$.
  Yen diwiwiti nganggo input kosong~$w = \emptyseq$,
  mesin iki ora bakal mandheg: mesin iki lumaku
  ing daur tanpa wates, nanging ora ngowahi pita
  utawa mindhah sirah. Kabeh konfigurasine padha
  (kahanan, posisi sirah, lan isi pita padha).
  Kita bisa netepake !!{structure} winates~$\Struct{M''}$
  kang nyukupi $!T(M,\emptyseq) \land !E(M,\emptyseq)$
  (gladhen). Nanging, kita bisa ngowahi~$!T(M,w)$
  kanthi cara kang trep supaya !!{structure}s
  kaya mangkono ora mungkin.

  \begin{prob}
    Tetepna $M$ minangka mesin Turing kang mung nduwé
    kahanan $q_0$ lan instruksi tunggal
    $\delta(q_0,\TMblank) = \tuple{q_0,\TMblank,\TMstay}$.
    Cathetan koreksi sumber OLPL-162: kahanan tujuan saiki dibenerake dadi q-nol kang siji-sijine, lan assignment penerus ing siji dibenerake supaya kabeh kanggo struktur M-prime-prime.
    Tetepna $\Domain{M''} = \{0, 1, 2\}$,
    $\Assign{\prime}{M''}(0) =
    \Assign{\prime}{M''}(1) = 1$ lan
    $\Assign{\prime}{M''}(2) = 2$, lan
    $\Assign{<}{M''} = \{\tuple{0,1}, \tuple{1,1}, \tuple{2,2}\}$.
    Tetepna $\Assign{\Obj Q_{q_0}}{M''}$,
    $\Assign{\Obj S_{\TMblank}}{M''}$, lan
    $\Assign{\Obj S_{\TMendtape}}{M''}$ saengga
    $!T(M,\emptyseq)$ lan $!E(M,\emptyseq)$
    dadi bener, banjur terangna alesane. Pituduh:
    Gatekna yen $\delta(q_0, \TMendtape)$ ora ditegesi.
    Priksa supaya
    \begin{align*}
    & \Obj Q_{q_0}(\num{1}, \num{n}) \land \Obj S_{\TMendtape}(\num{0}, \num{n})
      \land 
      \lforall[x][(\num{0} < x \lif \Obj S_\TMblank(x, \num{n}))] \text{\quad kanggo kabeh $n \in \Nat$}\\
    & \lexists[y][(\Obj Q_{q_0}(\num{0}, y) \land \Obj S_{\TMendtape}(\num{0}, y))]
    \end{align*}
    loro-lorone bener ing~$\Struct{M''}$.
  \end{prob}
\end{explain}

Gatekna !!{sentence}s kang njlentrehake tumindake
mesin Turing~$M$ nganggo input
$w = \sigma_{i_1}\dots\sigma_{i_k}$:
\begin{enumerate}
\item Aksioma kang njlentrehake angka lan~$<$
  (kaya ing definisi~$!T(M,w)$ ing \olref[rep]{sec}).
\item Aksioma kang njlentrehake konfigurasi input:
  kaya ing definisi~$!T(M,w)$.
\item Aksioma kang njlentrehake transisi saka
  sawijining konfigurasi menyang konfigurasi sabanjure:

Kanggo ing ngisor iki, tetepna $!A(x, y)$ kaya
sadurunge, lan tetepna
\[
  !B(y) \ident \lforall[x][(x < y \lif \eq/[x][y])].
\]
\begin{enumerate}
\item \ollabel{rep-right} Kanggo saben instruksi
  $\delta(q_i, \sigma) = \tuple{q_j, \sigma', \TMright}$,
  !!{sentence} iki:
\begin{align*}
& \lforall[x][\lforall[y][(
   (\Obj Q_{q_i}(x, y) \land \Obj S_{\sigma}(x, y)) \lif {}]] \\
&\qquad   (\Obj Q_{q_j}(x', y') \land \Obj S_{\sigma'}(x, y') \land
!A(x, y) \land !B(y')))
\end{align*}
\item \ollabel{rep-left} Kanggo saben instruksi
  $\delta(q_i, \sigma) = \tuple{q_j, \sigma', \TMleft}$,
  !!{sentence} iki
\begin{align*}
& \lforall[x][\lforall[y][
    ((\Obj Q_{q_i}(x', y) \land \Obj S_{\sigma}(x', y)) \lif {}]]\\
& \qquad   (\Obj Q_{q_j}(x, y') \land \Obj S_{\sigma'}(x', y') \land
!A(x', y) \land !B(y'))) \land {}\\
& \lforall[y][((\Obj Q_{q_i}(\num{0}, y) \land \Obj S_{\sigma}(\num{0},
    y)) \lif {}]\\
& \qquad (\Obj Q_{q_j}(\num{0}, y') \land \Obj S_{\sigma'}(\num{0},
  y') \land !A(\num{0}, y) \land !B(y')))
\end{align*}
Cathetan koreksi sumber OLPL-163 lan SOL6-F058: cabang obah kiwa saka kothak saliyane nol saiki nganggo pangreksan kanggo kothak asal kang ditulisi lan syarat B kanggo wektu sabanjure, kaya transisi liyane.
\item \ollabel{rep-stay} Kanggo saben instruksi
  $\delta(q_i, \sigma) = \tuple{q_j, \sigma', \TMstay}$,
  !!{sentence} iki:
\begin{align*}
& \lforall[x][\lforall[y][(
   (\Obj Q_{q_i}(x, y) \land \Obj S_{\sigma}(x, y)) \lif {}]] \\
&\qquad (\Obj Q_{q_j}(x, y') \land \Obj S_{\sigma'}(x, y') \land
!A(x, y) \land !B(y')))
\end{align*}
\end{enumerate}
Kaya kang katon, !!{sentence}s kang njlentrehake
transisine~$M$ padha karo !!{sentence} kang cocog
ing~$!T(M,w)$, kajaba kita nambahake $!B(y')$
ing pungkasan. $!B(y')$~njamin yen angka $y'$
tumrap konfigurasi ``sabanjure'' béda saka kabeh
angka sadurunge $\Obj 0$, $\Obj 0'$, \dots.
% \item Sentences that express consistency conditions:
% \begin{enumerate}
%   \item\ollabel{rep-con-symb}%
%   !!^a{sentence} that says that no square can have more than one
%   symbol on it at any given time:
%   \[\lforall[x][\lforall[y][(\Obj S_{\sigma}(x, y) \lif \lnot \Obj
%   S_{\sigma'}(x, y))]],\]
%   for every pair $\sigma \neq \sigma'$ of tape symbols of~$T$.
%   \item\ollabel{rep-con-state}%
%   !!^a{sentence} that says that $M$ cannot be in more than one state
%   at any given time:
%   \[\lforall[x][\lforall[y][(\Obj Q_{q}(x, y) \lif
%   \lforall[z][\lnot \Obj
%   Q_{q'}(z, y))]],\]
%   for every pair $q \neq q'$ of states of~$T$.
%   \item\ollabel{rep-con-tape}%
%   !!^a{sentence} that says that $M$ cannot be on more than one tape
%   square at any given time:
%   \[\lforall[x][\lforall[y][(\Obj Q_{q}(x, y) \lif
%   \lforall[z][\eq/[z][y]\lif \lnot \Obj
%   Q_{q'}(z, y))]],\]
%   for every pair $q$, $q'$ of states of~$T$.
% \end{enumerate}
\end{enumerate}
Tetepna $!T'(M, w)$ minangka konjungsi kabeh
!!{sentence}s ing ndhuwur kanggo mesin Turing~$M$
lan input~$w$.

\begin{lem}\ollabel{lem:halts-sat}
  Yen $M$ kang diwiwiti nganggo input~$w$ mandheg,
  mula $!T'(M,w) \land !E(M,w)$ nduwé model winates.
\end{lem}

\begin{proof}
  Tetepna $\Struct{M'}$ kaya ing bukti
  \olref[ver]{lem:halt-if-valid}, kajaba
  \begin{align*}
    \Domain{M'} & = \{0, \dots, n\},\\
    \Assign{\prime}{M'}(x) & = 
    \begin{cases}
      x + 1 &\text{yen $x < n$}\\
      n &\text{yen ora,}
    \end{cases} \\
    \tuple{x,y} \in \Assign{<}{M'} &\text{yen lan mung yen $x < y$ utawa $x = y = n$,}
  \end{align*}
  ing ngendi $n = \max(k+2,\len{w}+1)$ lan $k$~iku
  angka paling cilik saengga $M$ kang diwiwiti
  nganggo input~$w$ wis mandheg sawisé~$k$
  langkah. Pamriksan yen
  % Cathetan wates model winates diwenehake ing prosa sawise pratelan model.
  $\Sat{M'}{!T'(M,w) \land !E(M,w)}$
  kita tinggalake dadi gladhen. Cathetan koreksi sumber OLPL-161: wates saiki luwih gedhe tinimbang wektu mandheg, kabeh posisi kang ditekani lan dawa input. Kanthi posisi sirah wiwit siji, kothak kang ditekani paling tengen ora ngluwihi k kang ditambah siji. Transisi nganti mandheg ora tekan pucuk domain kang duwe daur. Cathetan OLPL-164: tandha formula E saiki dicocogake.
\end{proof}

\begin{prob}
  Rampungna bukti \olref[tur][und][tra]{lem:halts-sat}
  kanthi mbuktekake yen
  % Cathetan konsistensi simbol lema diwenehake sawise rumus gladhen.
  $\Sat{M'}{!T'(M,w) \land !E(M,w)}$. Cathetan koreksi sumber OLPL-164: formula ing gladhen saiki T-prime lan E kanthi tandha formula, padha karo lemma kang dibuktekake.
\end{prob}

\begin{lem}\ollabel{lem:sat-halts}
  Yen $!T'(M,w) \land !E(M,w)$ nduwé model winates,
  mula $M$ kang diwiwiti nganggo input~$w$ mandheg.
\end{lem}

\begin{proof}
  Kita nuduhake kontraposisine. Upamane $M$ kang
  diwiwiti nganggo~$w$ ora mandheg. Yen
  $!T'(M,w) \land !E(M,w)$ ora nduwé model babar pisan,
  bukti wis rampung. Mula upamane $\Struct{M}$
  % Cathetan asumsi model kanggo bukti kontraposisi diwenehake ing prosa.
  model saka~$!T'(M,w) \land !E(M, w)$. Cathetan koreksi sumber OLPL-165: asumsi model saiki T-prime lan E, mula syarat B bisa didudut saka saben transisi.
  Kita kudu nuduhake yen model iki ora mungkin winates.
  
  Kita bisa mbuktekake, kaya ing
  \olref[ver]{lem:config}, yen~$M$ kang diwiwiti
  nganggo input~$w$ durung mandheg sawisé~$n$
  langkah kanthi n positif, mula $!T'(M,w) \Entails !C(M, w, n)
  \land !B(\num{n})$. Amarga $M$ kang diwiwiti
  nganggo input~$w$ ora mandheg, mula
  $!T'(M,w) \Entails !C(M, w, n) \land
  !B(\num{n})$ kanggo kabeh~$n \in \Nat$ kang positif. Cathetan koreksi sumber SOL6-F064: ora ana aksioma B ing wektu nol. B ing siji saka transisi kapisan, banjur induksi menehi B ing saben wektu positif; iki cukup kanggo bukti domaine tanpa wates.
  Gatekna yen miturut \olref[rep]{prop:mlessk},
  $!T'(M,w) \Entails \num{k} < \num{n}$
  kanggo kabeh~$k < n$. Kajaba iku,
  $!B(\num{n}) \Entails \num{k} < \num{n} \lif
  \eq/[\num{k}][\num{n}]$. Dadi,
  $\Sat{M}{\eq/[\num{k}][\num{n}]}$ kanggo
  kabeh~$k < n$. Tegese, tanpa wates akeh
  suku~$\num{k}$ mau kudu nduwé nilai kang
  béda-béda ing~$\Struct{M}$. Nanging iki
  mbutuhake $\Domain{M}$ tanpa wates, mula
  $\Struct{M}$ ora bisa dadi model winates
  saka~$!T'(M,w) \land !E(M, w)$.
\end{proof}

\begin{prob}
  Rampungna bukti \olref[tur][und][tra]{lem:sat-halts}
  kanthi mbuktekake yen~$M$ kang diwiwiti
  nganggo input~$w$ durung mandheg sawisé~$n$
  langkah kanthi n positif, mula $!T'(M,w) \Entails !B(\num{n})$.
\end{prob}

\begin{thm}[Teorema Trakhtenbrot]
  \ollabel{thm:trakhtenbrodt}
  Ora bisa diputusake apa sembarang !!{sentence}
  logika tataran kapisan nduwé model winates
  (yaiku bisa dicukupi sacara winates).
\end{thm}

\begin{proof}
  Upamane ana mesin Turing~$F$ kang bisa mutusake
  masalah panyukupan sacara winates. Banjur yen
  diwenehi sembarang mesin Turing~$M$ lan
  input~$w$, kita bisa ngitung
  ukara~$!T'(M,w) \land !E(M,w)$, banjur
  nggunakake~$F$ kanggo mutusake apa ukara iku
  nduwé model winates. Miturut
  \cref{tur:und:tra:lem:halts-sat,tur:und:tra:lem:sat-halts},
  ukara iku nduwé model winates yen lan mung yen
  $M$ kang diwiwiti nganggo input~$w$ mandheg. Cathetan OLPL-161/163/773: pangreksan kothak asal, syarat B ing kabeh transisi lan wates model winates saiki dibenerake kaya kang dijlentrehake ing ndhuwur.
  Mula kita bisa nggunakake $F$ kanggo ngrampungake
  masalah mandheg, kang wis dimangerteni ora bisa
  dirampungake.
\end{proof}

\begin{cor}\ollabel{cor:fproof-incomp}%
  Ora mungkin ana sistem !!{derivation} efektif kang ngemu
  jaminan semantis lan jangkep tumrap kasahihan
  \emph{winates}, yaiku
  sistem !!a{derivation} kang nduwé $\Proves !B$
  yen lan mung yen $\Sat{M}{!B}$ kanggo saben
  !!{structure} winates~$\Struct{M}$. Cathetan koreksi sumber SOL6-F066: efektif ateges bukti winates bisa dipriksa lan dienumerasi kanthi algoritma, kang dibutuhake ing reduksi sabanjure.
\end{cor}

\begin{proof}
  Gladhen.
\end{proof}

\begin{prob}
  Buktekna \olref[tur][und][tra]{cor:fproof-incomp}.
  Gatekna yen $!B$ dicukupi ing saben
  !!{structure} winates yen lan mung yen
  $\lnot !B$ ora bisa dicukupi sacara winates.
  Terangna sebabé panyukupan sacara winates
  iku semi-bisa diputusake miturut teges ing
  \olref[tur][und][uns]{thm:valid-ce}. Gunakna
  iki kanggo nuduhake yen seandhane ana
  sistem !!a{derivation} efektif tumrap kasahihan winates,
  panyukupan sacara winates bakal bisa diputusake.
\end{prob}


\end{document}
